\section{Программная реализация}
\label{sec:implementation}

\subsection{Расчётные модули}

Расчётная часть программы реализована на C++ и включает модули геометрии и сетки, теплофизических свойств, аппроксимации потоков и решения прямой и обратной задач. Общие методы решения линейных, нелинейных и оптимизационных задач отделены от модели осадочного бассейна, поэтому при подборе параметров используется тот же прямой решатель. GMRES и метод POUNDERS предоставляются библиотекой PETSc и её пакетом оптимизации TAO~\cite{petscManual,petsc1997}, а прямое LU-разложение выполняется библиотекой UMFPACK~\cite{umfpack2004}.

Средства на Python обеспечивают чтение входных таблиц, подготовку данных, запись результатов в XLSX и визуализацию. Для каждой конфигурации температура и отражательная способность определяются на исходной геометрии на кровле, подошве и в среднем по мощности слоя, что сохраняет единый формат результатов при изменении расчётной сетки. Для обратной задачи дополнительно выводятся найденные узловые значения базального потока.

\subsection{Представление новых слоёв и перенос полей}

Число ячеек нового слоя определяется по его мощности в конечной конфигурации и сохраняется в течение эпохи. До отложения материала нулевая высота заменяется минимальной положительной величиной $\varepsilon=0{,}05$~м, которая обеспечивает ненулевые расстояния между центрами и гранями при вычислении потоков. Для $N$~сжатых ячеек суммарная мощность составляет $N\varepsilon$, и её влияние на результат оценивается в разделе~\ref{sec:verification}.

По мере осадконакопления высоты ячеек увеличиваются снизу вверх в соответствии с зависимостью~\eqref{eq:material-position}. На рис.~\ref{fig:growth} этот процесс показан для двух первых конфигураций Модели~1 в сечении $x=1400$~м. Новый слой постепенно накапливается над существующим, а нижний слой опускается с сохранением своей мощности.

\reportfigure{04_growth.png} {Постепенное осадконакопление по двум первым конфигурациям Модели~1 в сечении $x=1400$~м. (а) начало эпохи при возрасте 200~млн~лет; (б), (в) промежуточные состояния; (г) конец эпохи при возрасте 187,5~млн~лет. Новый слой накапливается сверху, а исходный слой опускается с сохранением мощности. Температура определяется в центрах ячеек, а отражательная способность витринита в центрах горизонтальных границ ячеек.} {fig:growth}

При изменении сетки на границе эпох температура в новом центре ячейки определяется линейной интерполяцией сначала по глубине в двух соседних колонках, затем по горизонтали. Аррениусовские интегралы интерполируются в пределах каждого сохраняющегося слоя, чтобы учитывать температурную историю соответствующего материала. Ячейкам ещё не отложившегося слоя задаётся температура поверхности, а интегралы вычисляются с индивидуальных моментов отложения $t_{\mathrm{dep}}$.

\subsection{Упрощение исходной геометрии}

Разбиение сетки на колонки зависит от узлов исходных границ, поэтому для уменьшения их числа используется алгоритм Рамера--Дугласа--Пекера (RDP)~\cite{ramer1972,douglas1973}, который заменяет участок ломаной отрезком, если отклонение промежуточных узлов не превышает заданного допуска. При большем отклонении участок разделяется в точке его максимума и та же процедура повторяется для двух полученных частей.

Горизонтальная ось сетки формируется с учётом границ всех слоёв. При упрощении сохраняются узлы, задающие области существования слоёв, их положительную мощность и выклинивание. На рис.~\ref{fig:rdp} приведены три участка одной границы Модели~2 при разных допусках и один участок сетки до и после упрощения.

\reportfigure{05_rdp.png} {Упрощение исходной геометрии Модели~2 алгоритмом RDP. (а)--(в) три участка одной границы с узлами исходной ломаной и после упрощения с допусками 1 и 5~м. Линии смещены по вертикали для наглядности. В легендах указаны допуск и число точек на показанном участке. (г), (д) один участок расчётной сетки без упрощения и с допуском 5~м.} {fig:rdp}

Для Модели~2 при допуске 5~м число ячеек уменьшается с 120335 до 3943, а время прямого расчёта сокращается с 205,41 до 4,95~с (табл.~\ref{tab:rdp}). По сравнению с аналогичным расчётом без упрощения геометрии относительные евклидовы нормы различий средних по слоям температуры и отражательной способности составляют 1,31 и 1,12\% соответственно. Максимальное локальное различие отражательной способности $E_{R_o}^{\max}$ достигает 6,91\%. Таким образом, алгоритм в данном примере позволил уменьшить время расчёта в 41~раз ценой небольшой потери точности.

\begin{table}[htbp]
\centering\small
\caption{Модель~2: упрощение геометрии и относительные отличия решения.}\label{tab:rdp}
\begin{tabular}{lrrrrr}
\toprule
Допуск, м & Ячеек & Время, с & $E_T$, \% & $E_{\Ro}$, \% & Макс. Ro, \% \\
\midrule
0 & 120335 & 205,41 & 0,00 & 0,00 & 0,00 \\
1 & 5288 & 7,79 & 0,56 & 0,92 & 3,48 \\
5 & 3943 & 4,95 & 1,31 & 1,12 & 6,91 \\
10 & 3723 & 3,62 & 1,32 & 1,18 & 6,90 \\
25 & 2955 & 2,72 & 1,36 & 1,11 & 7,49 \\
\bottomrule
\end{tabular}
\end{table}


\subsection{Расчёт отражательной способности витринита}

Аррениусовские интегралы $I_i$, определяющие отражательную способность витринита, сохраняются в центрах горизонтальных границ ячеек, и для колонки из $N$~ячеек используется $N+1$~точка расчёта интегралов, тогда как температурное поле задаётся в центрах ячеек. Такое размещение позволяет учитывать разные моменты отложения на верхней и нижней границах одной ячейки без введения дополнительных температурных неизвестных, а температура на внутренней границе определяется по значениям в соседних центрах с учётом тепловых сопротивлений в виде

\begin{equation}
 T_\sigma = \frac{r_N T_P + r_P T_N}{r_P + r_N},\qquad
 r_P = \frac{d_{P\sigma}}{\lambda_P},\quad
 r_N = \frac{d_{N\sigma}}{\lambda_N}.
\end{equation}

На поверхности используется заданная температура, а на нижней границе осадочного чехла её значение определяется экстраполяцией из центра нижней ячейки по базальному потоку и теплопроводности. Для каждой точки расчёта интегралов сохраняется момент отложения материала, который используется вместе с температурами на концах временного шага при вычислении приращений интегралов.

При линейном изменении температуры материальной точки во времени приращение~\eqref{eq:ro-increment} вычисляется по значениям первообразной на концах интервала, причём её явный вид

\begin{equation}
 \Phi_i(T) = T e^{-\beta_i/T} + \beta_i\operatorname{Ei}(-\beta_i/T),
\end{equation}
показывает, что непосредственный расчёт требует обращения к интегральной экспоненте для каждой реакции и каждой точки, поэтому для уменьшения этих затрат используется табулирование. После перехода к безразмерному аргументу $\chi=\beta_i/T$ из первообразной выделяются экспоненциальный и степенной множители, а оставшаяся медленно меняющаяся часть обозначается через $H(\chi)$ и сохраняется в таблице
\begin{gather}
 H(\chi) = \chi^2 e^\chi\left(\frac{e^{-\chi}}{\chi} + \operatorname{Ei}(-\chi)\right),\\
 \Phi_i(T) = \frac{T^2}{\beta_i}e^{-\beta_i/T} H(\beta_i/T).
\end{gather}

По результатам численных экспериментов выбрана таблица из 8192~значений $H$ в равномерных узлах интервала $20\leq\chi\leq200$, которая строится однократно и используется для всех реакций и точек расчёта. Равномерное расположение узлов позволяет определять интервал интерполяции по целой части нормированной координаты аргумента без поиска по таблице, а линейная интерполяция заменяет повторные обращения к интегральной экспоненте.

Для проверки точности табулирования рассматривается охлаждение однородного неподвижного слоя мощностью 4000~м, на поверхности которого задана постоянная температура, а на нижней границе~--- нулевой тепловой поток. В этой задаче температурная история известна из аналитического решения одномерного уравнения теплопроводности, которое имеет вид
\begin{equation}
 T_*(z,t) = 283{,}15 + 200\exp\left(-\frac{\lambda k^2}{C_0}t\right)\sin(kz),
 \qquad k = \frac{\pi}{8000\,\text{м}},
\end{equation}
что позволяет оценить погрешность вычисления отражательной способности независимо от погрешности численного решения задачи теплопереноса.

При теплопроводности $\lambda=2$~Вт/(м$\cdot$К) и объёмной теплоёмкости $C_0=3\cdot10^6$~Дж/(м$^3\cdot$К) сравнение проводится в 81~точке по глубине на интервале длительностью 1~млн~лет, разбитом на 200~временных шагов. Эталонные значения $R_o$ получены квадратурой по непрерывной температурной истории, причём удвоение числа интервалов квадратуры изменяет результат менее чем на $3{,}2\cdot10^{-7}\%$.

При линейном представлении температуры внутри временного шага максимальная относительная ошибка отражательной способности составляет 0,00590\%, тогда как дополнительное различие между табулированным и непосредственным вычислением первообразной для той же линейной истории не превышает $1{,}4\cdot10^{-8}\%$. Следовательно, в рассмотренной задаче погрешность определяется преимущественно приближением температурной истории внутри шага, а влияние табулирования существенно меньше (рис.~\ref{fig:easy-speed},~б).

\reportfigure{06_easy_ro.png} {Сравнение непосредственного вычисления первообразной и табулирования при расчёте отражательной способности витринита. (а) время вычисления аррениусовских интегралов для 65536~точек и 20~реакций; (б) относительная ошибка $R_o$ по глубине в задаче охлаждения однородного слоя. Эталон для панели (б) получен квадратурой по аналитической температурной истории.} {fig:easy-speed}

Для 65536~точек и 20~реакций табулирование сокращает время вычисления аррениусовских интегралов в 6,55~раза по сравнению с непосредственным вычислением первообразной (рис.~\ref{fig:easy-speed},~а), а доля этого этапа в полном времени программы приведена в табл.~\ref{tab:stages}.

\subsection{Повторные расчёты и выбор метода решения линейных систем}

При подборе базального потока геометрическая история остаётся заданной, поэтому геометрические данные, описания сеток эпох и коэффициенты интерполяции в точки наблюдений подготавливаются один раз для всех прямых расчётов. Начальная стационарная температура вычисляется для каждого нового потока, а запись промежуточных полей при оптимизации отключена.

Для выбора линейного решателя сопоставлены GMRES(100)/ILU(2) и LU-разложение в библиотеке UMFPACK~\cite{umfpack2004} на трёх моделях. Измерения выполнены на процессоре Intel Core i7-12700H с одним вычислительным потоком, компилятором GCC~15.2 с оптимизацией и библиотекой PETSc~3.24.4. Использованы шаги $\Delta x=200$~м, $\Delta z=50$~м, $\Delta t=1$~млн~лет, минимальная высота $\varepsilon=0{,}05$~м, допуск RDP~5~м и относительные допуски линейного и нелинейного решений $10^{-8}$. В табл.~\ref{tab:solvers} время каждого этапа приведено в секундах и процентах от соответствующего времени UMFPACK, без подготовки данных и записи XLSX средствами Python.

GMRES сокращает время решения линейных систем в 1,44--4,05~раза и полного прямого расчёта в 1,31--2,55~раза, при этом максимальное различие температурных полей между методами не превышает $3\cdot10^{-11}\,{}^\circ$C. По распределению времени в табл.~\ref{tab:stages} на линейное решение приходится 49,5--71,4\%, а на отражательную способность~--- 2,3--4,1\%, поэтому выбранный итерационный метод ускоряет наиболее затратный этап при практически совпадающих результатах.

\begin{table}[htbp]
\centering\small
\renewcommand{\arraystretch}{1.12}
\caption{Время расчёта в секундах и процентах от времени UMFPACK для того же этапа и модели.}\label{tab:solvers}
\begin{tabular}{ll D{,}{,}{2.3} D{,}{,}{3.1} D{,}{,}{2.3} D{,}{,}{3.1}}
\toprule
\multirow{2}{*}{Модель} & \multirow{2}{*}{Решатель} & \multicolumn{2}{c}{Линейное решение} & \multicolumn{2}{c}{Прямой расчёт} \\
\cmidrule(lr){3-4}\cmidrule(lr){5-6}
\multicolumn{1}{c}{} & \multicolumn{1}{c}{} & \multicolumn{1}{c}{с} & \multicolumn{1}{c}{\%} & \multicolumn{1}{c}{с} & \multicolumn{1}{c}{\%} \\
\midrule
\multirow{2}{*}{Модель~1} & GMRES/ILU(2) & 1,160 & 24,7 & 2,291 & 39,2 \\
 & UMFPACK & 4,698 & 100,0 & 5,839 & 100,0 \\
\addlinespace[3pt]
\multirow{2}{*}{Модель~2} & GMRES/ILU(2) & 2,471 & 27,3 & 4,995 & 43,1 \\
 & UMFPACK & 9,062 & 100,0 & 11,586 & 100,0 \\
\addlinespace[3pt]
\multirow{2}{*}{Модель~3} & GMRES/ILU(2) & 16,049 & 69,2 & 22,486 & 76,2 \\
 & UMFPACK & 23,178 & 100,0 & 29,511 & 100,0 \\
\bottomrule
\end{tabular}
\end{table}

\begin{table}[htbp]
\centering\small
\renewcommand{\arraystretch}{1.12}
\caption{Время этапов и доля полного прямого расчёта.}\label{tab:stages}
\begin{tabular}{l D{,}{,}{2.3} D{,}{,}{3.1} D{,}{,}{2.3} D{,}{,}{3.1} D{,}{,}{2.3} D{,}{,}{3.1} }
\toprule
\multirow{2}{*}{Этап} & \multicolumn{2}{c}{Модель~1} & \multicolumn{2}{c}{Модель~2} & \multicolumn{2}{c}{Модель~3} \\
\cmidrule(lr){2-3}\cmidrule(lr){4-5}\cmidrule(lr){6-7}
\multicolumn{1}{c}{} & \multicolumn{1}{c}{с} & \multicolumn{1}{c}{\%} & \multicolumn{1}{c}{с} & \multicolumn{1}{c}{\%} & \multicolumn{1}{c}{с} & \multicolumn{1}{c}{\%} \\
\midrule
Стационарный старт & 0,002 & 0,1 & 0,002 & 0,0 & 0,010 & 0,0 \\
Геометрия и сетка & 0,027 & 1,2 & 0,050 & 1,0 & 0,170 & 0,8 \\
Перенос & 0,009 & 0,4 & 0,014 & 0,3 & 0,043 & 0,2 \\
Сборка & 0,929 & 40,5 & 1,751 & 35,1 & 5,235 & 23,3 \\
Линейное решение & 1,160 & 50,6 & 2,471 & 49,5 & 16,049 & 71,4 \\
EASY\%Ro & 0,093 & 4,1 & 0,165 & 3,3 & 0,528 & 2,3 \\
Вывод & 0,068 & 3,0 & 0,541 & 10,8 & 0,423 & 1,9 \\
Прочее & 0,004 & 0,2 & 0,001 & 0,0 & 0,028 & 0,1 \\
\midrule
Всего & 2,291 & 100,0 & 4,995 & 100,0 & 22,486 & 100,0 \\
\bottomrule
\end{tabular}
\end{table}


Таким образом, программная реализация сочетает вывод полей на исходной геометрии с ускорением расчёта за счёт упрощения границ, табулирования и выбора линейного решателя. Однократная подготовка геометрических данных дополнительно сокращает повторяющиеся операции при восстановлении базального потока.
